\section{Introduction}

\begin{figure}[htbp]
    \centering
    \includegraphics[width=\linewidth]{iclr2026/figures/Figure1-genund_teaser_refined_v13_vector.pdf}
\caption{
\textbf{When does generation help understanding?}
(a) Increasing I2I supervision can steadily improve downstream I2T performance, while the effect depends strongly on the training recipe.
(b) Across visual capabilities, transfer is highly task-dependent: some I2I--I2T pairs yield large gains, whereas others interfere with understanding (blue/red denote positive/negative transfer; bubble size denotes magnitude).
(c) This selectivity is also reflected inside the model, where gradient alignment between I2I and I2T objectives varies sharply across parameter groups and is strongest in pre-attention normalization.
Together, these observations motivate our study of \emph{when, where, and why} generation helps understanding.
}
    \label{fig:teaser}
\end{figure}
% und -> gen, gen-> und?; prev think doesn't help
% 1. start with data motivation; prev paper showed that generation can't really help, intuitively doesn;t make sense, we challenge this; thinking with images work can be explained here, thinking trace can help, co training cannot help
When does image generation help visual understanding?
Prior work has shown that visual understanding can improve image generation, yet whether training on generation can in turn improve understanding remains unclear~\citep{tong2024metamorph,kang2026transferability}.
This has led to a view that understanding helps generation, while generation provides substantially less benefit to understanding~\citep{tong2024metamorph, ICLR2026_2aadf45d, NEURIPS2025_59d4e18a, li2026manzano}.
However, this asymmetry is counterintuitive: image generation provides dense supervision over object appearance, spatial relationships, and geometry, which are also essential for visual understanding tasks such as recognition, counting, spatial reasoning, and 3D perception.
Indeed, recent evidence suggests that generative pretraining can learn general-purpose visual representations useful for a broad range of downstream vision tasks~\citep{visionbanana2026}.
If such supervision can transfer, generation data could provide an additional source of training signal for understanding tasks.
We therefore revisit this assumption and systematically study when generation supervision improves visual understanding.

In this work, we study generation as a source of training supervision for visual understanding.
This is orthogonal to recent ``thinking with images'' methods, which improve visual reasoning by generating intermediate images at inference time~\citep{hu2024visual,gu2025thinkmorph,qin2025chain,fu2025refocus,li2025zebra,li2025imagine}.
We instead ask whether generation can improve understanding without requiring image generation at inference time.
Unified multimodal models provide a natural setting for studying this transfer, as generation and understanding are learned within the same model~\citep{metaqueries,wu2025harmonizing,bagel,han2025tv2tv}.
We systematically study (1) whether such transfer occurs and how it depends on the training recipe and amount of generation supervision, (2) which generation tasks benefit which understanding capabilities, and (3) what explains successful transfer.

We first ask: \emph{Can generation supervision improve visual understanding, and under what training setup does this transfer emerge?}
To isolate generation-to-understanding transfer, we construct paired generation and understanding tasks that require the same underlying visual operation while differing only in output modality.
As illustrated in~\autoref{fig:scaling}(a), the two tasks share the same input and underlying problem, but express the answer in different forms.
For example, given a shuffled image, the I2I task reconstructs the correctly ordered image, while the corresponding I2T task predicts the patch order in text.
After training with the generation objective, we evaluate the model on its
corresponding I2T task.
We compare several training recipes for combining generation and understanding data, including mixed I2I-I2T training and two-stage training with I2I training followed by I2T training.
We find that the two-stage recipe yields the strongest and most consistent gains.
Under this recipe, I2I supervision consistently improves the paired I2T tasks, and the gains continue to increase as more I2I data is added, as shown in~\autoref{fig:teaser}(a).

We next ask: \emph{Which kinds of generation data help which kinds of understanding tasks?}
Prior work has studied transfer relationships among visual tasks~\citep{zamir2018taskonomy},
but answering this question across generation and understanding requires a common capability space that spans both output modalities.
We therefore introduce \methodname, a unified capability taxonomy that organizes I2I supervision tasks and I2T capabilities within the same visual capability space.
\methodname is organized around the three foundational problems of computer vision (3R)~\citep{malik2016three}: \textit{Recognition}, \textit{Reconstruction}, and \textit{Reorganization}. 
Within each family, we derive fine-grained understanding capabilities bottom-up from existing benchmarks, and organize I2I objectives according to the visual information they supervise.
This shared structure allows us to systematically measure transfer from each I2I task to each understanding capability.
Applying the best training recipe across \methodname reveals a highly structured generation-to-understanding transfer map (Fig.~\ref{fig:teaser}(b)).
Some gains follow intuitive capability correspondences, such as Z-depth improving metric 3D reasoning and object pointing improving counting. 
Others are less obvious: 2D keypoint supervision substantially improves both lighting understanding and multi-view reasoning. 
Moreover, the same generation task can help some capabilities while interfering with others. 

% last talk about gradients analysis, and how this guides experiments (with results)
% \xd{how about start with: If transfer arises from compatible representation learning, successful I2I tasks should induce parameter updates aligned with those from the corresponding I2T tasks.}
Finally, we ask: \emph{What makes generation supervision transfer successfully to a particular understanding task?}
If transfer arises from compatible representation learning, successful
I2I-I2T pairs should induce parameter updates aligned across the two objectives.
This hypothesis is motivated by prior work showing that conflicting task gradients
can impair multi-task optimization and that gradient-based interactions can be
used to estimate task affinity~\citep{yu2020gradient,fifty2021efficient}.
Therefore, we analyze this through gradient alignment between the I2I and I2T objectives.
We find that stronger downstream transfer is associated with higher gradient alignment, with the strongest alignment concentrated in normalization parameters, particularly pre-attention normalization, as shown in~\autoref{fig:teaser}(c).
These results suggest that optimization compatibility between generation and understanding objectives may be one factor associated with successful transfer.

% finally summarize:
Our contributions are summarized as follows:
(1) We introduce a controlled framework for studying training-time transfer from image generation to visual understanding, and show that I2I supervision can substantially improve paired I2T tasks under the right training recipe, with gains increasing as more generation data is added.
(2) We introduce \methodname, a unified capability taxonomy for I2I and I2T tasks, and use it to reveal highly structured, task-dependent generation-to-understanding transfer.
(3) We analyze gradient alignment in the controlled setting and find that stronger transfer is associated with more aligned optimization signals, particularly in normalization parameters.

% comments from baifeng shi
% 1. start with data motivation; prev paper showed that generation can't really help, intuitively doesn;t make sense, we challenge this; thinking with images work can be explained here, thinking trace can help, co training cannot help
% 2. cite think with images; conclude -> separate, the questions we have asked in this paper; we want to study co-training can help or not; recipe; taskonomy, ...
% 3. bagel is too early, maybe later 
% 4. delete also

% in teaser figure only keep the matrix